\begin{afterword}
\summaryhead

The post continues ``Update Yourself Incrementally'' from the day before. It suggests that when people meet a new argument against a belief, they avoid lowering their confidence by recalling the supporting arguments they already knew.

The example is hypothetical. Freedonia suspects its neighbour Sylvania of causing meteor strikes, on three pieces of evidence. Each time someone offers a counterargument, such as the trade between the countries or Sylvania's lack of a space program, the Freedonian repeats the three known points, judges the balance ``3 to 1,'' and ends up more confident than before.

The author calls this double-counting: the old arguments were already included in the old confidence. Counting all the data twice would be bad enough; counting only the favourable part twice is a farce. Even when the balance still favours a belief, each new piece of contrary evidence must lower its probability. It appears to the author that people do this when they meet a new counterargument; the author must keep constant vigilance against doing it too.

\responsehead

The post makes one point, and it is correct. Evidence already counted cannot be counted again to offset new evidence. The Freedonia story shows the error plainly: the counterarguments are serious ones, trade worth billions and the lack of a space program, and each is dismissed by the same three points. The comparison with a scientist who counts every subject twice makes it clear why the practice misleads.

What the post does not show is that the error is common. That people behave this way is put first as a suggestion (``Here I suggest'') and later as an impression (``it does appear to me''). The only case is hypothetical, and the only other support is the author's report of having to guard against it. The claim may be true, but the post does not give the reader a way to judge.

The post follows from the previous day's post, which said that contrary evidence must lower a belief even when the belief survives; this one names a way people avoid doing so, and applies it to the author as well as to the reader.

In short: a correct and clearly shown point about double-counting evidence, with the claim that people commonly do it supported only by the author's impression.
\end{afterword}
